\documentclass[10pt,a4paper]{amsart}
\usepackage{amsmath,amssymb,amsthm,graphicx}
% Minimal retained preamble for lane-s5's SIGMA 2025-010 card, taken from
% sources/sigma-2025-010-fermionic-basis-in-conformal-field-theory-an/pinned-original-source.tex
% (doi:10.3842/SIGMA.2025.010, SIGMA 21 (2025), article 010, CC BY 4.0).
%
% Retained from the source preamble: the equation numbering, the theorem declarations
% under the shared counter, and the macro lines the retained spans actually use.
%
% Omitted and why: \pdfoutput=1 and the \ifpdf / \documentclass[pdftex]{sigma} block
% (pdfTeX-only primitive and the journal class, which texlive does not carry; per the
% lane-x1 precedent the assembler supplies the article/amsart baseline class), and the
% SIGMA front-matter commands, which belong to the journal layout and are not used by
% the retained spans.
\usepackage{amsmath,amssymb,amsthm,mathrsfs,mathtools,xfrac}

\numberwithin{equation}{section}

\newtheorem{Theorem}{Theorem}[section]
\newtheorem{Proposition}[Theorem]{Proposition}
{ \theoremstyle{definition}
\newtheorem{Definition}[Theorem]{Definition}
\newtheorem{Remark}[Theorem]{Remark} }

% macro lines copied verbatim from the source preamble (lines 19-29 of the pinned file)
\newcommand\at[2]{\left.#1\right|_{#2}}
\DeclareMathOperator{\tr}{Tr}
\DeclareMathOperator{\res}{res}
\newcommand{\vac}{{\rm vac}}
\newcommand{\qbinom}[3]{\big[
                        \begin{smallmatrix}
                        #1\\ #2
                        \end{smallmatrix}
                        \big]_#3}

\newcommand{\mathcomment}[1]{\Big\slash\text{#1}\Big\slash}
\title{A new representation of the function Psi in terms of the basic fermionic operators}
\author{Adler, Sergei, Boos, Hermann}
\date{doi:10.3842/SIGMA.2025.010; CC BY 4.0}
\begin{document}
\maketitle
\setcounter{section}{3}
\setcounter{subsection}{0}
\setcounter{Theorem}{0}
\setcounter{equation}{0}
\noindent\textit{Editorial note.} All delivered mathematical text is the author's own, copied from the pinned SIGMA article source (source\_id doi:10.3842/SIGMA.2025.010); no statement or proof step was rewritten, reordered or removed.\par\medskip
\noindent\textit{Editorial note.} The proof of this proposition is not adjacent to its statement: the paper proves it in the appendix section 'Proof of Proposition prop:Psi', and the proof span is that appendix section, from its heading down to the sentence that finishes the proof (CONTROL 8.8.157: a far-placed named proof is delivered as the single proof span).\par\medskip
\noindent\textit{Editorial note.} The source uses no proof environment at all, so the span guard's proof-environment check is skipped for this source; the delivered proof span is the author's own appendix section.\par\medskip
\noindent\textit{Editorial note.} The counted spans carry no bibliography citation; the paper's references section is outside both spans.\par\medskip
\noindent\textit{Editorial note.} The preamble span points at the top of the pinned article source (its own header, lines 1-33, up to the line before the begin-document), used verbatim as staged.\par\medskip
\noindent\textit{Editorial note.} The card carries a minimal card-local preamble instead of the article's own header: the header uses the pdfTeX-only \textbackslash{}pdfoutput primitive and the journal class sigma.cls, neither of which the host's TeX Live provides. The retained lines (equation numbering, theorem declarations, macro definitions) are copied verbatim from the pinned source header.\par\medskip
\noindent\textit{Allowed standard prerequisites (editorial).}\begin{itemize}
\item The setup of the fermionic basis and the known integral representation of Psi\_nu. Same paper, Section 2 (Review of the previous results), in force in the statement. Fixes the operators, the kernel F-bar and the equation the proof compares against.
\item The epsilon-expansion of the functions zeta\_0, zeta\_R and of the kernel F-bar. Same paper, the expansions displayed at proof lines 723-740. Supplies the order-by-order comparison used in the induction.
\item Residue calculus for the fermionic operators. Same paper, Section 2, the residue conventions. Underlies every displayed computation in the appendix.
\end{itemize}\medskip
\begin{align}
& \Psi(l, \kappa)+\frac{{\rm i}f\kappa}{l\bigl(l+\frac{{\rm i}}{2\nu}\bigr)}=
-\sum_{r=1}^{k} \res_h\biggl[\frac{{\rm e}^{-\sfrac{h x_r^+(\kappa)}{f\kappa}}-{\rm e}^{-\sfrac{h x_r^-(\kappa)}{f\kappa}}}{h(l+h)}S(h)\biggr] \nonumber\\
&\qquad{}-\frac{2{\rm i}}{f\kappa} \sum\limits_{n=0}^{\infty}\frac{1}{n!} \int_0^\infty \frac{{\rm d}x}{2\pi}
\bigg\{
\res_h\biggl[\frac{{\rm e}^{-\sfrac{hx}{f\kappa}}}{l+h}S(h)\biggr] \bar{F}(x, \kappa)^n \biggl(-\frac{\partial}{\partial x}\biggr)^n
\bigg\}_{{\rm even}}
\log\bigl(1+{\rm e}^{-2\pi x}\bigr).\!\!\!\label{1_Psi_1}
\end{align}

\begin{align}
\Psi_\nu(s, p):=\frac{1}{2\nu {\rm i}}\Psi\Bigl(\frac{s}{2\nu {\rm i}}, p\Bigr), \qquad
S_\nu(h):=S\biggl(\frac{h}{2\nu {\rm i}}\biggr).\label{1_Rescale}
\end{align}

\begin{align}\label{1_Zetas}
&\zeta_0(s, p)=\zeta_0(s, p), \\
&\zeta_1(s, p)=\res_h\bigl[\sigma_1(h)\zeta_0(h, p) p^{-1}(1+s-h) \zeta_0(1+s-h, p)\bigr], \cr
&\zeta_2(s, p)=\res_h\bigl[\bigl( \sigma_2(h)\zeta_0(h, p) + \sigma_1(h)\zeta_1(h, p) \bigr) p^{-1}(1+s-h) \zeta_0(1+s-h, p) \bigr] \cr
&              \qquad{}+\frac{1}{2}\res_h\res_{h'}\bigl[\sigma_1(h) \zeta_0(h, p)\sigma_1\bigl(h'\bigr)\zeta_0\bigl(h', p\bigr)p^{-2}\bigl(1+s-h-h'\bigr)_2\zeta_0\bigl(2+s-h-h', p\bigr)\bigr]. \nonumber
\end{align}

\begin{Proposition}
\label{prop:Psi}
The function $\Psi_\nu(s, p)$ is given by the formula
\begin{align}\label{1_Psi_Proposition}
&\Psi_\nu(s, p)|0\rangle=\exp\{\res_h[(S_\nu(h)-1)\Psi_\nu(h, p)\Phi_+(h)]\} \mathcal{F}(s, p, \Phi_-)|0\rangle,
\end{align}
where $\mathcal{F}(s, p, \Phi_-)$ is a functional of fields $\Phi_{-}$,
\begin{align*}
&\mathcal{F}(s, p, \Phi_-)\\
&\qquad{}  =\sum_{n=0}^{\infty}\frac{p^{-n}}{n!}\res_{h_1}\cdots\res_{h_n} \Phi_-(h_1)\cdots\Phi_-(h_n) \Bigl(1+s-\sum h_j\Bigr)_n \zeta_0\Bigl(n+s-\sum h_j, p\Bigr).
\end{align*}
\end{Proposition}

\section{Proof of Proposition \ref{prop:Psi}}\label{appendix:PsiProof}
In this appendix, we give the proof of Proposition \ref{prop:Psi}. For simplicity, we consider the case without any particle-hole excitations. We want to show that the equation~\eqref{1_Psi_Proposition} gives the same result as the equation~\eqref{1_Psi_1} in every order of the $\epsilon$-expansion. In ``renormalised'' variables~\eqref{1_Rescale}, the representation~\eqref{1_Psi_1} of $\Psi_\nu$ takes the form
\begin{align}
&\Psi_\nu(s, p)-p\Psi_0(s, p) \nonumber\\
&\qquad{}=-\frac{2}{p} \sum_{n=0}^{\infty} \frac{1}{n!} \int_{0}^{\infty} \frac{{\rm d}x}{2\pi}
        \biggl\{\res_h\biggl[ \frac{{\rm e}^{\frac{{\rm i}hx}{p}}}{s+h} S(h) \biggr] \bar{F}(x, p)^n \biggl(-\frac{\partial}{\partial x}\biggr)^n \biggr\}_{\text{even}}\log\bigl(1+{\rm e}^{-2\pi x}\bigr),\label{A_Psi_HGS}
\end{align}
where the function $\bar{F}(x, p)$ is given by
\begin{align}\label{A_1}
\bar{F}(x, p)=x+{\rm i}\res_h\big[ {\rm e}^{\frac{{\rm i}hx}{p}} S(h)\Psi_\nu(h, p) \big].
\end{align}
Let us introduce the expansion of the function $\bar{F}(x, p)$ with respect to $\epsilon$,
\begin{align*}
\bar{F}(x, p)=\sum_{j=0}^{\infty} \bar{F}_j(x, p) \epsilon^j.
\end{align*}
From~\eqref{A_1}, it follows that
\begin{align}
&\bar{F}_0(x, p)=x+{\rm i}\res_h\bigl[ {\rm e}^{\frac{{\rm i}hx}{p}} \zeta_0(h, p) \bigr] = x+\res_h\biggl[ {\rm i} {\rm e}^{\frac{{\rm i}hx}{p}} \frac{p}{h(h-1)} \biggr] = {\rm i}p \sum_{m=2}^{\infty} \frac{1}{m!} \biggl(\frac{{\rm i}x}{p}\biggr)^m, \nonumber\\
&\bar{F}_1(x, p)={\rm i}\res_h\bigl[ {\rm e}^{\frac{{\rm i}hx}{p}} \sigma_1(h)\zeta_0(h, p) \bigr], \cr
&\bar{F}_2(x, p)={\rm i}\res_h\bigl[ {\rm e}^{\frac{{\rm i}hx}{p}} (\sigma_2(h)\zeta_0(h, p) + \sigma_1(h)\zeta_1(h, p)) \bigr], \qquad \dots.
\label{A_F_Epsilon}
\end{align}
The $\epsilon$-expansion of the $\Psi(h, p)$ can be obtained from~\eqref{A_Psi_HGS},
\begin{gather}
\Psi_\nu(s, p)-p\Psi_0(s, p)\nonumber\\
\qquad{}=-\frac{2}{p} \sum_{n=0}^{\infty} \frac{1}{n!} \int_{0}^{\infty} \frac{{\rm d}x}{2\pi}
        \bigg\{ \res_h\biggl[ \frac{{\rm e}^{\frac{{\rm i}hx}{p}}}{s+h} \biggr]
        \bigg(F_0(x, p)^n
              \epsilon\cdot nF_0(x, p)^{n-1}F_1(x, p)  \cr
\qquad\hphantom{=}{} +\epsilon^2 \biggl(nF_0(x, p)^{n-1}F_2(x, p) + \frac{n(n-1)}{2}F_0(x, p)^{n-1}F_1(x, p)^2\biggr) + \cdots \bigg) \! \biggl(-\frac{\partial}{\partial x}\biggr)^n \bigg\}_{\text{even}} \cr
\qquad\hphantom{=+}{}\times \log\bigl(1+{\rm e}^{-2\pi x}\bigr).\label{A_HGS_Epsilon}
\end{gather}

As it can be seen, in every order of the $\epsilon$-expansion of the function $\Psi(s, p)$ one has to calculate the integral of the form
\begin{align}\label{A_Int}
-\frac{2}{p} \sum_{n=0}^{\infty} \frac{1}{n!} \int_{0}^{\infty} \frac{{\rm d}x}{2\pi} \biggl\{ {\rm e}^{-\frac{{\rm i}sx}{p}} \bar{F}_0(x, p)^n \biggl(-\frac{\partial}{\partial x}\biggr)^{n+n'} \biggr\}_{\text{even}} \log\bigl(1+{\rm e}^{-2\pi x}\bigr)
\end{align}
for some $n'\geq 0$. To do so, let us introduce the following notation:
\begin{align*}%\label{A_u}
\frac{1}{n!}\Biggl[ \sum_{k=2}^{\infty} \frac{x^k}{k!} \Biggr]^n=
\sum_{m=2n}^{\infty} u_{n, m} x^m.
\end{align*}
The numbers $u_{n, m}$ can be written down explicitly in the form
\begin{align*}
u_{n, m}=\frac{1}{n!} \sum_{\substack{2\leq k_1 \leq\cdots\leq k_n \leq m-2(n-1), \\ k_1+\cdots+k_n=m}} \frac{1}{k_1! \cdots k_n!}.
\end{align*}
The most important property of $u_{n, m}$ relevant to this proof is the following sum:
\begin{align}\label{A_u_1}
\sum_{\substack{0 \leq n \leq r, \\ 0 \leq k \leq n+r}} (-1)^{n+k} \frac{(n+r)!}{k!} u_{n, r+n-k} s^k=(-1)^r (s+1)_r.
\end{align}
Then the left-hand side of the formula~\eqref{A_Int} can be written as
\begin{align}
&-\frac{2}{p} \sum_{n=0}^{\infty} \frac{1}{n!} \int_{0}^{\infty} \frac{{\rm d}x}{2\pi} \biggl\{ {\rm e}^{-\frac{{\rm i}sx}{p}} \bar{F}_0(x, p)^n \biggl(-\frac{\partial}{\partial x}\biggr)^{n+n'} \biggr\}_{\text{even}} \log\bigl(1+{\rm e}^{-2\pi x}\bigr)\nonumber\\
&\qquad{}=-\frac{2}{p} \sum_{\substack{n \geq 0, \\ k \geq 0, \\ m \geq 2n}} \int_{0}^{\infty} \frac{{\rm d}x}{2\pi}
    \biggl\{{\rm i}^{-k+n+m} p^{n-m-k} x^{k+m} \frac{s^k}{k!} u_{n, m} \biggl(-\frac{\partial}{\partial x}\biggr)^{n+n'} \biggr\}_{\text{even}} \log\bigl(1+{\rm e}^{-2\pi x}\bigr)\cr
&\qquad{} =-2\sum_{\substack{n \geq 0,\\ k \geq 0, \\ m \geq 2n}} {\rm i}^{-k+n+m} p^{n-m-k-1} \frac{s^k}{k!} u_{n, m} (m+k)! \bigl(1-2^{-m-k-1+n+n'}\bigr)\nonumber\\
&\qquad\hphantom{=-2\sum_{\substack{n \geq 0,\\ k \geq 0, \\ m \geq 2n}}}{}
\times\frac{\zeta\bigl(m+k-n-n'+2\bigr)}{(2\pi)^{m+k-n-n'+2}}.\label{A_2}
\end{align}
In the final line, the equality
\begin{align*}
\int_{0}^{\infty} \frac{{\rm d}x}{2\pi} x^{m} \biggl(-\frac{\partial}{\partial x}\biggr)^n \log\bigl(1+{\rm e}^{-2\pi x}\bigr) = m!\bigl(1-2^{-m-1+n}\bigr) \frac{\zeta(m-n+2)}{(2\pi)^{m-n+2}}
\end{align*}
was used and it was assumed that $k+m-n-n'$ is even.

Let us denote $m+k-n=r$. Since $u_{n, m}=0$ for $m<2n$, the sums over $k$ and $n$ in~\eqref{A_2} are actually finite. Using the property~\eqref{A_u_1} of the numbers $u_{n, m}$, we continue the derivation
\begin{align}
\eqref{A_2}&{}=-2\sum_{\substack{n \geq 0,\\ k \geq 0,\\ r \geq 0}} p^{-r-1} (-1)^{n+k} {\rm i}^{r} \frac{(r+n)!}{k!} u_{n, r-k+n} s^k \bigl(1-2^{-r+n'-1}\bigr) \frac{\zeta\bigl(r-n'+2\bigr)}{(2\pi)^{r-n'+2}}\nonumber\\
           &{}=-2\sum_{\substack{0 \leq n \leq r,\\ 0 \leq k \leq r-n, \\ r \geq 0}} p^{-r-1} (-1)^{n+k} {\rm i}^{r} \frac{(r+n)!}{k!} u_{n, r-k+n} s^k \bigl(1-2^{-r+n'-1}\bigr) \frac{\zeta\bigl(r-n'+2\bigr)}{(2\pi)^{r-n'+2}}\cr
           &{}=-2\sum_{r \geq 0} (-{\rm i})^{r} p^{-r-1} (s+1)_r \bigl(1-2^{-r-n'+1}\bigr)\frac{\zeta\bigl(r-n'+2\bigr)}{(2\pi)^{r-n'+2}}.\label{A_3}
\end{align}
We can denote $r-n'+2=2m$ since it is even. Then
\begin{align*}
\eqref{A_3}&=-2\sum_{m=-n'+2}^{\infty} (-{\rm i})^{2m+n'-2} p^{-2m+1-n'} (s+1)_{2m+n'-2} \bigl(1-2^{-2m+1}\bigr)\frac{\zeta(2m)}{(2\pi)^{2m}}\\
                     &=-2({\rm i}p)^{-n'}(1+s)_{n'}\frac{1}{s+n'}\sum_{m=1}^{\infty} p^{-2m+1} \bigl(s+n'\bigr)_{2m-1} \bigl(1-2^{-2m+1}\bigr) (-{\rm i})^{2m-2} \frac{|B_{2m}|}{2}\cr
                     &\quad{} -2({\rm i}p)^{-n'} \sum_{m=-n'+2}^{0} {\rm i}^{2m-2} p^{-2m+1} (1+s)_{2m+n'-2} \bigl(1-2^{-2m+1}\bigr)\frac{\zeta(2m)}{(2\pi)^{2m}} \cr
                     &=({\rm i}p)^{-n'} (1+s)_{n'} \zeta_{R}\bigl(s+n', p\bigr)-{\rm i}({\rm i}p)^{-n'+1}(1+s)_{n'-2}\cr
                     &\quad{} -2({\rm i}p)^{-n'} \sum_{m=-n'+2}^{0} {\rm i}^{2m-2} p^{-2m+1} (1+s)_{2m+n'-2} \bigl(1-2^{-2m+1}\bigr)\frac{\zeta(2m)}{(2\pi)^{2m}} .\nonumber
\end{align*}
For $n'\geq2$, the only term that contributes to the sum in the last line is the one with $m=0$. Otherwise the sum is zero. Therefore, we get to the following result:
\begin{align}
&-\frac{2}{p} \sum_{n=0}^{\infty} \frac{1}{n!} \int_{0}^{\infty} \frac{{\rm d}x}{2\pi} \biggl\{ {\rm e}^{-\frac{{\rm i}sx}{p}} \bar{F}_0(x, p)^n \biggl(-\frac{\partial}{\partial x}\biggr)^{n+n'} \biggr\}_{\text{even}} \log\bigl(1+{\rm e}^{-2\pi x}\bigr) \cr
&\qquad =
\begin{cases}
({\rm i}p)^{-n'} (1+s)_{n'} \zeta_{R}\bigl(s+n', p\bigr) & \text{for } n'=0, 1, \\
({\rm i}p)^{-n'} (1+s)_{n'} \zeta_{0}\bigl(s+n', p\bigr) & \text{for } n'\geq 2.
\end{cases}\label{A_Integral_Answer}
\end{align}

From this result for $n'=0$, it follows immediately that the $\epsilon^0$-term of the expansion~\eqref{A_Psi_HGS} is~$\zeta_0(s, p)$ which is in agreement with Proposition \ref{prop:Psi}. Now let us turn our attention to the first order with respect to $\epsilon$. From~\eqref{A_HGS_Epsilon} and~\eqref{A_F_Epsilon}, we get
\begin{align*}
\zeta_1(s, p)
&=-\frac{2}{p}\Biggl( \sum_{n=1}^{\infty} \frac{{\rm i}}{n!} \int_{0}^{\infty} \frac{{\rm d}x}{2\pi}
\biggl\{
    {\rm e}^{-\frac{{\rm i}sx}{p}} n \bar{F}_0(x, p)^{n-1} \res_h\biggl[ {\rm e}^{\frac{{\rm i}hx}{p}} \sigma_1(h)\zeta_0(h, p) \biggr]
    \biggl(-\frac{\partial}{\partial x}\biggr)^{n}
\biggr\}_{\text{even}}\cr
&\quad{}+\sum_{n=0}^{\infty} \frac{1}{n!} \int_{0}^{\infty} \frac{{\rm d}x}{2\pi}
\biggl\{
    \res_h\biggl[\frac{{\rm e}^{\frac{{\rm i}hx}{p}}}{s+h}\sigma_1(h)\biggr] \bar{F}_0(x, p)^n
    \biggl(-\frac{\partial}{\partial x}\biggr)^{n}
\biggr\}_{\text{even}}\Biggr) \log\bigl(1+{\rm e}^{-2\pi x}\bigr).
\end{align*}
We can shift the summation index in the first sum and then use~\eqref{A_Integral_Answer} to calculate both sums. This gives us
\begin{align}\label{A_Zeta1_1}
\zeta_1(s, p)=\res_h\big[ \sigma_1(h)\zeta_0(h, p) p^{-1} (1+s-h) \zeta_R(1+s-h, p) \big] + \res_h\biggl[ \frac{\sigma_1(h)\zeta_R(-h)}{s+h} \biggr].
\end{align}
To bring the answer to the final form we can use the duality property of $S(h)$ which in the new variables simply reads
\begin{align}
S(h)S(-h)=1 \quad\Rightarrow\quad &\sigma_1(h)+\sigma_1(-h)=0, \nonumber\\
                             &\sigma_1(h)\sigma_1(-h)+\sigma_2(h)+\sigma_2(-h)=0, \cr
                             &\ldots.\label{A_S_Duality}
\end{align}
Then the last term in the formula~\eqref{A_Zeta1_1} can be modified and added to the first one,
\begin{align*}
\zeta_1(s, p)&{}=\res_h\bigl[ \sigma_1(h)\zeta_0(h, p) p^{-1} (1+s-h) \zeta_R(1+s-h, p) \bigr] + \res_h\biggl[ \frac{\sigma_1(h)\zeta_R(h)}{s-h} \biggr]  \\
&{}=\res_h\bigl[ \sigma_1(h)\zeta_0(h, p) p^{-1} (1+s-h) \zeta_R(1+s-h, p) \bigr] + \res_h\biggl[ \frac{\sigma_1(h)\zeta_0(h)}{s-h} \biggr]\\
&{}=\res_h\bigl[ \sigma_1(h)\zeta_0(h, p) p^{-1} (1+s-h) \zeta_0(1+s-h, p) \bigr].
\end{align*}
That is exactly the result~\eqref{1_Zetas} we obtain from the exponential representation.

To proceed with higher orders of the expansion of $\Psi_\nu(s, p)$, let us introduce a shorthand notation for the integral~\eqref{A_Psi_HGS} in form of the ``contraction''
\begin{align*}%\label{A_Contraction_Def}
&\langle g, f \rangle_{n'} \\
&\quad{}=-\frac{2}{p} \sum_{n=0}^{\infty} \frac{1}{n!} \int_{0}^{\infty} \frac{{\rm d}x}{2\pi}
    \biggl\{
        \res_h\biggl[{\rm e}^{\frac{{\rm i}hx}{p}} \frac{f(h)}{s+h} \bar{F}_0(x, p)^n f(x, p)\biggr]
        \biggl(-\frac{\partial}{\partial x}\biggr)^{n+n'}
    \biggr\}_{\text{even}} \!\log\bigl(1+{\rm e}^{-2\pi x}\bigr). \nonumber
\end{align*}
Then the formula~\eqref{A_Integral_Answer}, for example, takes a simple form
\begin{align*}
\langle 1, 1 \rangle_{n'} =
\begin{cases}
({\rm i}p)^{-n'} (1+s)_{n'} \zeta_{R}\bigl(s+n', p\bigr) & \text{for } n'=0, 1, \\
({\rm i}p)^{-n'} (1+s)_{n'} \zeta_{0}\bigl(s+n', p\bigr) & \text{for } n'\geq 2.
\end{cases}
\end{align*}
For $\zeta_1(s, p)$, we have
\begin{align*}
\zeta_1(s, p)=\bigl\langle 1, \bar{F}_1\bigr\rangle_1 + \langle \sigma_1, 1 \rangle_0.
\end{align*}
We apply~\eqref{A_S_Duality} to the last term and add to the first one to replace $\zeta_R(1+s, p)$ by $\zeta_0(1+s, p)$ as it was done above.

Now let us discuss the general case $\zeta_k(s, p)$. We will prove Proposition \ref{prop:Psi} by induction. Suppose it is correct for some $k-1$. Using the ``contraction'' defined above, we write
\begin{align}
\zeta_k(s, p)&= \bigl\langle 1, \bar{F}_{k} \bigr\rangle_1 + \bigl\langle 1, \bar{F}_{k-1} \bar{F}_1 \bigr\rangle_2 + \bigl\langle 1, \bar{F}_{k-2} \bar{F}_{2} \bigr\rangle_2 + \frac{1}{2!} \bigl\langle 1, \bar{F}_{k-2}\bar{F}_1^2\bigr\rangle + \cdots \nonumber\\
             &\quad{}+ \langle \sigma_k(h), 1 \rangle_0
              + \bigl\langle \sigma_{k-1}(h), \bar{F}_1 \bigr\rangle_1
              + \biggl( \bigl\langle \sigma_{k-2}, \bar{F}_2\bigr\rangle_1 + \frac{1}{2}\bigl\langle \sigma_{k-2}, \bar{F}_1^2 \bigr\rangle_2 \biggr)+ \cdots .\label{A_ZetaK_Gen}
\end{align}
The very first term is
\begin{align}
&\bigl\langle 1, \bar{F}_{k} \bigr\rangle_1 \nonumber\\
&\qquad{}= \res_h \bigl[ (\sigma_1(h)\zeta_{k-1}(h)+\cdots+\sigma_{k-1}(h)\zeta_1(h, p))(1+s-h) p^{-1} \zeta_R(1+s-h, p) \bigr].\label{A_1_FK}
\end{align}
If it had the function $\zeta_0(1+s-h, p)$ instead of $\zeta_R(1+s-h, p)$, the first line would have given us the exponent~\eqref{1_Psi_Proposition}, which is the direct consequence of the integral~\eqref{A_Integral_Answer}. It can be shown that the terms coming from the second line of~\eqref{A_ZetaK_Gen} are exactly the missing singular terms we need to add to~\eqref{A_1_FK} to replace $\zeta_R(1+s-h)$ by $\zeta_0(1+s-h, p)$. Indeed, let us apply the duality~\eqref{A_S_Duality},
\begin{align}
\langle \sigma_k, 1 \rangle_0 &{}= \res_h \biggl[ \frac{\sigma_k(h)}{s+h} \zeta_0(-h, p)\biggr] \nonumber\\
                                 &{}= \res_h \biggl[ \frac{\sigma_k(h)}{s-h} \zeta_0(h, p)\biggr] + \sum_{i=1}^{k-1} \res_h \biggl[ \frac{\sigma_i(h) \sigma_{k-i}(-h)}{s-h} \zeta_0(h, p)\biggr].\label{A_S_Duality_Gen}
\end{align}
Now in~\eqref{A_ZetaK_Gen} take the terms which are the ``contractions'' of $\sigma_i(h)$ with some powers of $\bar{F}_j$. Once again, we will get the function $\zeta_R$ in one of such ``contractions''
\begin{align*}
\begin{split}
&\bigl\langle \sigma_i, \bar{F}_{k-i} \bigr\rangle_1=\res_h\res_{h'} \biggl[ \frac{\sigma_i(h)}{s+h}\bigl(\sigma_{k-i}\bigl(h'\bigr)\zeta_0(h', p) + \cdots+ \sigma_1\bigl(h'\bigr)\zeta_{k-i-1}\bigl(h'\bigr)\bigr)\\
&\hphantom{\bigl\langle \sigma_i, \bar{F}_{k-i} \bigr\rangle_1=\res_h\res_{h'} \biggl[}{}
\times\bigl(1-h-h'\bigr)p^{-1} \zeta_R\bigl(1-h-h', p\bigr) \biggr].
\end{split}
\end{align*}
The missing singular term is exactly the term proportional to $\sigma_i(h)$ coming from~\eqref{A_S_Duality_Gen}. We write, therefore,
\begin{align*}
&\bigl\langle \sigma_i, \bar{F}_{k-i} \bigr\rangle_1 + \res_h \biggl[ \frac{\sigma_i(h) \sigma_{k-i}(-h)}{s-h} \zeta_0(h, p)\biggr] = \res_h\res_{h'} \biggl[ \frac{\sigma_i(h)}{s-h}\\
&\qquad{}\times\bigl(\sigma_{k-i}\bigl(h'\bigr)\zeta_0\bigl(h', p\bigr) + \cdots + \sigma_1\bigl(h'\bigr)\zeta_{k-i-1}\bigl(h'\bigr)\bigr)\bigl(1-h-h'\bigr)\frac{1}{p}\zeta_0\bigl(1-h-h', p\bigr) \biggr].
\end{align*}
Now, by the assumption of induction all terms the terms which are a ``contraction'' of $\sigma_i(h)$ with some powers of $\bar{F}_j$ sum up into
\begin{align*}
\sum_{i=1}^{k-1} \res_h\biggl[\frac{\sigma_i(h) \zeta_{k-i}(h, p)}{s-h} \biggr],
\end{align*}
and these are exactly the singular terms we need to be able to replace $\zeta_R$ by $\zeta_0$ in~\eqref{A_1_FK}. This observation finishes the proof of Proposition \ref{prop:Psi}.

\end{document}
